\FigureLayoutDeclare{img/2019/beijing_liberal/q8_fig1.png}{6.0cm}{11bp 15bp 14bp 11bp}

\chapter{2019年北京卷（文）}

\section{单选题}

\begin{problem}
已知集合 \(A=\{x\mid -1<x<2\}\)，\(B=\{x\mid x>1\}\)，则 \(A\cup B=\)（\quad）.

\choices
    {\(\left(-1,1\right)\)}
    {\(\left(1,2\right)\)}
    {\(\left(-1,+\infty\right)\)}
    {\(\left(1,+\infty\right)\)}
\end{problem}

\begin{answer}
C
\end{answer}

\begin{problem}
已知复数 \(z=2+\i\)，则 \(z\overline z=\)（\quad）.

\choices
    {\(\sqrt3\)}
    {\(\sqrt5\)}
    {\(3\)}
    {\(5\)}
\end{problem}

\begin{answer}
D
\end{answer}

\begin{problem}
下列函数中，在区间 \(\left(0,+\infty\right)\) 上单调递增的是（\quad）.

\choices
    {\(y=\sqrt{x}\)}
    {\(y=2^{-x}\)}
    {\(y=\log_{\frac12}x\)}
    {\(y=\dfrac1x\)}
\end{problem}

\begin{answer}
A
\end{answer}

\begin{problem}
执行如图所示的程序框图，输出的 \(s\) 值为（\quad）.

\bitmapfigure[max width=6.0cm]{img/2019/beijing_liberal/q4_fig1.png}

\choices
    {\(1\)}
    {\(2\)}
    {\(3\)}
    {\(4\)}
\end{problem}

\begin{answer}
B
\end{answer}

\begin{problem}
已知双曲线 \(\dfrac{x^2}{a^2}-y^2=1\)（\(a>0\)）的离心率是 \(\sqrt5\)，则 \(a=\)（\quad）.

\choices
    {\(6\)}
    {\(4\)}
    {\(2\)}
    {\(\dfrac12\)}
\end{problem}

\begin{answer}
D
\end{answer}

\begin{problem}
设函数 \(f(x)=\cos x+b\sin x\)（\(b\) 为常数），则“\(b=0\)”是“\(f(x)\) 为偶函数”的（\quad）.

\choices
    {充分而不必要条件}
    {必要而不充分条件}
    {充分必要条件}
    {既不充分也不必要条件}
\end{problem}

\begin{answer}
C
\end{answer}

\begin{problem}
在天文学中，天体的明暗程度可以用星等或亮度来描述. 两颗星的星等与亮度满足 \(m_2-m_1=\dfrac52\lg\dfrac{E_1}{E_2}\)，其中星等为 \(m_k\) 的星的亮度为 \(E_k\)（\(k=1\)，\(2\)）. 已知太阳的星等是 \(-26.7\)，天狼星的星等是 \(-1.45\)，则太阳与天狼星的亮度的比值为（\quad）.

\choices
    {\(10^{10.1}\)}
    {\(10.1\)}
    {\(\lg 10.1\)}
    {\(10^{-10.1}\)}
\end{problem}

\begin{answer}
A
\end{answer}

\begin{problem}
如图，\(A\)，\(B\) 是半径为 \(2\) 的圆周上的定点，\(P\) 为圆周上的动点，\(\angle APB\) 是锐角，大小为 \(\beta\). 图中阴影区域的面积的最大值为（\quad）.

\bitmapfigure[max width=6.0cm]{img/2019/beijing_liberal/q8_fig1.png}

\choices
    {\(4\beta+4\cos\beta\)}
    {\(4\beta+4\sin\beta\)}
    {\(2\beta+2\cos\beta\)}
    {\(2\beta+2\sin\beta\)}
\end{problem}

\begin{answer}
B
\end{answer}

\section{填空题}

\begin{problem}
已知向量 \(\bs a=(-4,3)\)，\(\bs b=(6,m)\)，且 \(\bs a\perp\bs b\)，则 \(m=\fillinblank{}\).
\end{problem}

\begin{answer}
\(8\)
\end{answer}

\begin{problem}
若 \(x\)，\(y\) 满足 \(x\le2\)，\(y\ge-1\)，\(4x-3y+1\ge0\)，则 \(y-x\) 的最小值为\fillinblank{}，最大值为\fillinblank{}.
\end{problem}

\begin{answer}
\(-3\)，\(1\)
\end{answer}

\begin{problem}
设抛物线 \(y^2=4x\) 的焦点为 \(F\)，准线为 \(l\). 则以 \(F\) 为圆心，且与 \(l\) 相切的圆的方程为\fillinblank{}.
\end{problem}

\begin{answer}
\((x-1)^2+y^2=4\)
\end{answer}

\begin{problem}
某几何体是由一个正方体去掉一个四棱柱所得，其三视图如图所示. 如果网格纸上小正方形的边长为 \(1\)，那么该几何体的体积为\fillinblank{}.

\bitmapfigure[max width=0.58\linewidth]{img/2019/beijing_liberal/q12_fig1.png}
\end{problem}

\begin{answer}
\(40\)
\end{answer}

\begin{problem}
已知 \(l\)，\(m\) 是平面 \(\alpha\) 外的两条不同直线. 给出下列三个论断：

\begin{circlelist}
\item \(l\perp m\)；

\item \(m\parallel\alpha\)；

\item \(l\perp \alpha\).
\end{circlelist}

以其中的两个论断作为条件，余下的一个论断作为结论，写出一个正确的命题：\fillinblank{}.
\end{problem}

\begin{answer}
若 \(l\perp m\)，\(l\perp\alpha\)，则 \(m\parallel\alpha\).
\end{answer}

\begin{problem}
李明自主创业，在网上经营一家水果店，销售的水果中有草莓，京白梨，西瓜，桃，价格依次为 \(60\) 元/盒，\(65\) 元/盒，\(80\) 元/盒，\(90\) 元/盒. 为增加销量，李明对这四种水果进行促销：一次购买水果的总价达到 \(120\) 元，顾客就少付 \(x\) 元. 每笔订单顾客网上支付成功后，李明会得到支付款的 \(80\%\).

\begin{enumerate}
\item 当 \(x=10\) 时，顾客一次购买草莓和西瓜各 \(1\) 盒，需要支付\fillinblank{} 元；

\item 在促销活动中，为保证李明每笔订单得到的金额均不低于促销前总价的七折，则 \(x\) 的最大值为\fillinblank{}.
\end{enumerate}
\end{problem}

\begin{answer}
\(130\)，\(15\)
\end{answer}

\section{解答题}

\begin{problem}
在 \(\triangle ABC\) 中，\(a=3\)，\(b-c=2\)，\(\cos B=-\dfrac12\).

\begin{enumerate}
\item 求 \(b\)，\(c\) 的值；

\item 求 \(\sin(B+C)\) 的值.
\end{enumerate}
\end{problem}

\begin{solution}
\begin{enumerate}
\item 由余弦定理，\(b^2=a^2+c^2-2ac\cos B\). 将 \(a=3\)，\(b=c+2\)，\(\cos B=-\dfrac12\) 代入，得
\[
(c+2)^2=3^2+c^2-2\times3\times c\times\left(-\frac12\right)
\]
整理得 \(c^2+4c+4=c^2+3c+9\)，解得 \(c=5\). 从而 \(b=c+2=7\).

\item 因为 \(\cos B=-\dfrac12\) 且 \(B\in(0,\pi)\)，所以 \(B=\dfrac{2\pi}{3}\)，\(\sin B=\dfrac{\sqrt3}{2}\).
在 \(\triangle ABC\) 中，由正弦定理 \(\dfrac{a}{\sin A}=\dfrac{b}{\sin B}\)，得
\[
\sin A=\frac{a\sin B}{b}=\frac{3\times\frac{\sqrt3}{2}}{7}=\frac{3\sqrt3}{14}
\]
又 \(A+B+C=\pi\)，所以 \(\sin(B+C)=\sin(\pi-A)=\sin A=\dfrac{3\sqrt3}{14}\).
\end{enumerate}
\end{solution}

\begin{problem}
设 \(\{a_n\}\) 是等差数列，\(a_1=-10\)，且 \(a_2+10\)，\(a_3+8\)，\(a_4+6\) 成等比数列.

\begin{enumerate}
\item 求 \(\{a_n\}\) 的通项公式；

\item 记 \(\{a_n\}\) 的前 \(n\) 项和为 \(S_n\)，求 \(S_n\) 的最小值.
\end{enumerate}
\end{problem}

\begin{solution}
\begin{enumerate}
\item 设等差数列 \(\{a_n\}\) 的公差为 \(d\). 因为 \(a_1=-10\)，所以 \(a_2=-10+d\)，\(a_3=-10+2d\)，\(a_4=-10+3d\).
由 \(a_2+10\)，\(a_3+8\)，\(a_4+6\) 成等比数列，得
\[
(a_3+8)^2=(a_2+10)(a_4+6)
\]
即 \((2d-2)^2=d(3d-4)\). 展开整理得 \(d^2-4d+4=0\)，即 \((d-2)^2=0\)，解得 \(d=2\).
故 \(\{a_n\}\) 的通项公式为 \(a_n=-10+2(n-1)=2n-12\).

\item 由（1）知 \(a_n=2n-12\). 当 \(n\le5\) 时，\(a_n<0\)；当 \(n=6\) 时，\(a_6=0\)；当 \(n\ge7\) 时，\(a_n>0\).
因此数列的前 \(n\) 项和 \(S_n\) 在 \(n=5\) 或 \(n=6\) 时取得最小值.
最小值为
\[
S_6=\frac{6(a_1+a_6)}{2}=\frac{6\times(-10+0)}{2}=-30
\]
\end{enumerate}
\end{solution}

\begin{problem}
改革开放以来，人们的支付方式发生了巨大转变. 近年来，移动支付已成为主要支付方式之一. 为了解某校学生上个月 \(A\)，\(B\) 两种移动支付方式的使用情况，从全校所有的 \(1000\) 名学生中随机抽取了 \(100\) 人，发现样本中 \(A\)，\(B\) 两种支付方式都不使用的有 \(5\) 人，样本中仅使用 \(A\) 和仅使用 \(B\) 的学生的支付金额分布情况如下：

\begin{center}
\begin{tabular}{|c|c|c|}
\hline
支付方式 & 不大于 \(2\,000\) 元 & 大于 \(2\,000\) 元 \\
\hline
仅使用 \(A\) & \(27\) 人 & \(3\) 人 \\
\hline
仅使用 \(B\) & \(24\) 人 & \(1\) 人 \\
\hline
\end{tabular}
\end{center}

\begin{enumerate}
\item 估计该校学生中上个月 \(A\)，\(B\) 两种支付方式都使用的人数；

\item 从样本仅使用 \(B\) 的学生中随机抽取 \(1\) 人，求该学生上个月支付金额大于 \(2\,000\) 元的概率；

\item 已知上个月样本学生的支付方式在本月没有变化. 现从样本仅使用 \(B\) 的学生中随机抽查 \(1\) 人，发现他本月的支付金额大于 \(2\,000\) 元. 结合（2）的结果，能否认为样本仅使用 \(B\) 的学生中本月支付金额大于 \(2\,000\) 元的人数有变化？说明理由.
\end{enumerate}
\end{problem}

\begin{solution}
\begin{enumerate}
\item 在抽取的 \(100\) 人样本中，两种支付方式都不使用的有 \(5\) 人，仅使用 \(A\) 的有 \(27+3=30\) 人，仅使用 \(B\) 的有 \(24+1=25\) 人. 因此样本中两种支付方式都使用的学生人数为
\[
100-5-30-25=40\text{ 人}
\]
样本中两种支付方式都使用的比例为 \(\dfrac{40}{100}=0.4\). 由此估计全校 \(1\,000\) 名学生中两种支付方式都使用的人数为 \(1\,000\times0.4=400\) 人.

\item 在样本中仅使用 \(B\) 的学生共有 \(24+1=25\) 人，其中支付金额大于 \(2\,000\) 元的有 \(1\) 人.
因此从样本仅使用 \(B\) 的学生中随机抽取 \(1\) 人，该学生上个月支付金额大于 \(2\,000\) 元的概率为
\[
P=\frac{1}{25}
\]

\item 能认为有变化. 理由：由（2）知，上个月样本仅使用 \(B\) 的学生中支付金额大于 \(2\,000\) 元的概率为 \(\dfrac{1}{25}=0.04\)，这是小概率事件. 在本月抽查 \(1\) 人即出现支付金额大于 \(2\,000\) 元，根据小概率事件在一次试验中几乎不可能发生的原理，有充分理由认为本月仅使用 \(B\) 的学生中支付金额大于 \(2\,000\) 元的人数发生了变化（增加）.（注：若说明单次抽样具有随机偶然性，认为不能确定有变化，言之成理亦可）
\end{enumerate}
\end{solution}

\begin{problem}
如图，在四棱锥 \(P-ABCD\) 中，\(PA\perp\) 平面 \(ABCD\)，底面 \(ABCD\) 为菱形，\(E\) 为 \(CD\) 的中点.

\begin{enumerate}
\item 求证：\(BD\perp\) 平面 \(PAC\)；

\item 若 \(\angle ABC=60^\circ\)，求证：平面 \(PAB\perp\) 平面 \(PAE\)；

\item 棱 \(PB\) 上是否存在点 \(F\)，使得 \(CF\parallel\) 平面 \(PAE\)？说明理由.
\end{enumerate}

\bitmapfigure[max width=0.52\linewidth]{img/2019/beijing_liberal/q18_fig1.png}
\end{problem}

\begin{solution}
\begin{enumerate}
\item 因为底面 \(ABCD\) 为菱形，所以 \(BD\perp AC\).
又因为 \(PA\perp\) 平面 \(ABCD\)，且 \(BD\subset\) 平面 \(ABCD\)，所以 \(PA\perp BD\)，即 \(BD\perp PA\).
又 \(AC\cap PA=A\)，且 \(AC\)，\(PA\subset\) 平面 \(PAC\)，所以 \(BD\perp\) 平面 \(PAC\).

\item 在菱形 \(ABCD\) 中，\(AD=CD\)，\(\angle ADC=\angle ABC=60^\circ\)，所以 \(\triangle ACD\) 为等边三角形.
因为 \(E\) 为 \(CD\) 的中点，所以 \(AE\perp CD\).
又 \(AB\parallel CD\)，所以 \(AE\perp AB\).
因为 \(PA\perp\) 平面 \(ABCD\)，且 \(AE\subset\) 平面 \(ABCD\)，所以 \(PA\perp AE\)，即 \(AE\perp PA\).
又 \(AB\cap PA=A\)，且 \(AB\)，\(PA\subset\) 平面 \(PAB\)，所以 \(AE\perp\) 平面 \(PAB\).
因为 \(AE\subset\) 平面 \(PAE\)，所以平面 \(PAB\perp\) 平面 \(PAE\).

\item 存在点 \(F\)，当 \(F\) 为棱 \(PB\) 的中点时，\(CF\parallel\) 平面 \(PAE\). 理由如下：
连接 \(AC\)，设 \(AC\cap BD=O\)，连接 \(OE\)，\(OF\).
在 \(\triangle BCD\) 中，\(O\) 为 \(BD\) 的中点，\(E\) 为 \(CD\) 的中点，所以 \(OE\parallel BC\) 且 \(OE=\dfrac12BC\).
因为 \(ABCD\) 为菱形，所以 \(AD\parallel BC\) 且 \(AD=BC\)，故 \(OE\parallel AD\) 且 \(OE=\dfrac12AD\).
取 \(PA\) 的中点为 \(M\)，连接 \(FM\)，\(ME\).
在 \(\triangle PAB\) 中，\(F\)，\(M\) 分别为 \(PB\)，\(PA\) 的中点，所以 \(FM\parallel AB\) 且 \(FM=\dfrac12AB\).
因为 \(AB\parallel CD\)，所以 \(FM\parallel CD\). 又 \(CE=\dfrac12CD\)，故 \(FM\parallel CE\) 且 \(FM=CE\)，因此四边形 \(CFME\) 为平行四边形，从而 \(CF\parallel ME\).
因为 \(ME\subset\) 平面 \(PAE\)，且 \(CF\not\subset\) 平面 \(PAE\)，所以 \(CF\parallel\) 平面 \(PAE\).
\end{enumerate}
\end{solution}

\begin{problem}
已知椭圆 \(C:\dfrac{x^2}{a^2}+\dfrac{y^2}{b^2}=1\) 的右焦点为 \((1,0)\)，且经过点 \(A(0,1)\).

\begin{enumerate}
\item 求椭圆 \(C\) 的方程；

\item 设 \(O\) 为原点，直线 \(l:y=kx+t\)（\(t\ne\pm1\)）与椭圆 \(C\) 交于两个不同点 \(P\)，\(Q\)，直线 \(AP\) 与 \(x\) 轴交于点 \(M\)，直线 \(AQ\) 与 \(x\) 轴交于点 \(N\)，若 \(\lvert OM\rvert\cdot\lvert ON\rvert=2\)，求证：直线 \(l\) 经过定点.
\end{enumerate}
\end{problem}

\begin{solution}
\begin{enumerate}
\item 因为椭圆 \(C\) 的右焦点为 \((1,0)\)，所以 \(c=1\).
又椭圆经过点 \(A(0,1)\)，代入方程得 \(\dfrac{0^2}{a^2}+\dfrac{1^2}{b^2}=1\)，解得 \(b^2=1\).
由 \(a^2=b^2+c^2=1+1=2\)，得椭圆 \(C\) 的方程为 \(\dfrac{x^2}{2}+y^2=1\).

\item 设 \(P(x_1,y_1)\)，\(Q(x_2,y_2)\). 将 \(y=kx+t\) 代入 \(\dfrac{x^2}{2}+y^2=1\)，得
\[
(1+2k^2)x^2+4ktx+2t^2-2=0
\]
判别式 \(\Delta=(4kt)^2-4(1+2k^2)(2t^2-2)=8(1+2k^2-t^2)>0\).
由韦达定理，\(x_1+x_2=-\dfrac{4kt}{1+2k^2}\)，\(x_1x_2=\dfrac{2(t^2-1)}{1+2k^2}\).
直线 \(AP\) 过点 \(A(0,1)\) 与 \(P(x_1,y_1)\)，其方程为 \(y-1=\dfrac{y_1-1}{x_1}x\). 令 \(y=0\)，得点 \(M\) 的横坐标为 \(x_M=\dfrac{x_1}{1-y_1}\). 同理，点 \(N\) 的横坐标为 \(x_N=\dfrac{x_2}{1-y_2}\).
从而 \(|OM|\cdot|ON|=|x_M x_N|=\left|\dfrac{x_1x_2}{(y_1-1)(y_2-1)}\right|\).
因为 \(P\)，\(Q\) 在直线 \(y=kx+t\) 上，所以
\[
(y_1-1)(y_2-1)=(kx_1+t-1)(kx_2+t-1)=k^2x_1x_2+k(t-1)(x_1+x_2)+(t-1)^2
\]
代入韦达定理化简得
\[
(y_1-1)(y_2-1)=\frac{(t-1)^2}{1+2k^2}
\]
因此
\[
\frac{x_1x_2}{(y_1-1)(y_2-1)}=\frac{\frac{2(t^2-1)}{1+2k^2}}{\frac{(t-1)^2}{1+2k^2}}=\frac{2(t+1)}{t-1}
\]
由 \(|OM|\cdot|ON|=2\)，得 \(\left|\dfrac{2(t+1)}{t-1}\right|=2\)，即 \(\left|\dfrac{t+1}{t-1}\right|=1\).
因为 \(t\ne\pm1\)，所以 \(t+1=-(t-1)\)，解得 \(t=0\).
此时 \(\Delta=8(1+2k^2)>0\) 恒成立，直线 \(l\) 的方程为 \(y=kx\)，恒过原点 \((0,0)\).
因此直线 \(l\) 经过定点 \((0,0)\).
\end{enumerate}
\end{solution}

\begin{problem}
已知函数 \(f(x)=\dfrac14x^3-x^2+x\).

\begin{enumerate}
\item 求曲线 \(y=f(x)\) 的斜率为 \(1\) 的切线方程；

\item 当 \(x\in\left[-2,4\right]\) 时，求证：\(x-6\le f(x)\le x\)；

\item 设 \(F(x)=\lvert f(x)-(x+a)\rvert\)（\(a\in\R\)），记 \(F(x)\) 在区间 \(\left[-2,4\right]\) 上的最大值为 \(M(a)\)，当 \(M(a)\) 最小时，求 \(a\) 的值.
\end{enumerate}
\end{problem}

\begin{solution}
\begin{enumerate}
\item 函数的导数为 \(f'(x)=\dfrac34x^2-2x+1\). 令 \(f'(x)=1\)，得 \(\dfrac34x^2-2x=0\)，解得 \(x_1=0\)，\(x_2=\dfrac83\).
当 \(x=0\) 时，\(f(0)=0\)，切线方程为 \(y=x\)；
当 \(x=\dfrac83\) 时，\(f\left(\dfrac83\right)=\dfrac{8}{27}\)，切线方程为 \(y-\dfrac{8}{27}=x-\dfrac83\)，即 \(y=x-\dfrac{64}{27}\).
所以曲线 \(y=f(x)\) 的斜率为 \(1\) 的切线方程为 \(y=x\) 或 \(y=x-\dfrac{64}{27}\).

\item 设 \(h(x)=f(x)-x=\dfrac14x^3-x^2=\dfrac14x^2(x-4)\).
当 \(x\in[-2,4]\) 时，\(x^2\ge0\)，\(x-4\le0\)，所以 \(h(x)\le0\)，即 \(f(x)\le x\).
又设 \(g(x)=f(x)-(x-6)=\dfrac14x^3-x^2+6\)，其导数为 \(g'(x)=\dfrac34x^2-2x=x\left(\dfrac34x-2\right)\).
当 \(x\in[-2,0)\) 时 \(g'(x)>0\)，\(g(x)\) 单调递增；当 \(x\in\left(0,\dfrac83\right)\) 时 \(g'(x)<0\)，\(g(x)\) 单调递减；当 \(x\in\left(\dfrac83,4\right]\) 时 \(g'(x)>0\)，\(g(x)\) 单调递增.
因为 \(g(-2)=0\)，\(g\left(\dfrac83\right)=\dfrac{98}{27}>0\)，所以当 \(x\in[-2,4]\) 时 \(g(x)\ge g(-2)=0\)，即 \(f(x)\ge x-6\).
综上，当 \(x\in[-2,4]\) 时，\(x-6\le f(x)\le x\).

\item 设 \(h(x)=f(x)-x=\dfrac14x^3-x^2\)，则 \(F(x)=|h(x)-a|\).
由（2）知，当 \(x\in[-2,4]\) 时，\(h'(x)=x\left(\dfrac34x-2\right)\)，极值点为 \(x=0\)，\(x=\dfrac83\).
计算端点与极值点处的函数值：\(h(-2)=-6\)，\(h(0)=0\)，\(h\left(\dfrac83\right)=-\dfrac{64}{27}\)，\(h(4)=0\).
所以 \(h(x)\) 在 \([-2,4]\) 上的最大值为 \(0\)，最小值为 \(-6\)，即对任意 \(x\in[-2,4]\)，有 \(-6\le h(x)\le0\).
当 \(a=-3\) 时，\(-3\le h(x)+3\le3\)，故 \(M(-3)=\max_{x\in[-2,4]}|h(x)+3|=3\).
若 \(a>-3\)，则取 \(x=-2\)，有 \(F(-2)=|-6-a|=a+6>3\)；
若 \(a<-3\)，则取 \(x=0\)，有 \(F(0)=|0-a|=-a>3\).
所以当且仅当 \(a=-3\) 时，\(M(a)\) 取得最小值 \(3\).
故当 \(M(a)\) 最小时，\(a=-3\).
\end{enumerate}
\end{solution}
